\documentclass[11pt]{article}
\usepackage[margin=1in]{geometry}
\usepackage{amsmath,amssymb,amsthm}
\usepackage{graphicx}
\usepackage[hidelinks]{hyperref}

\newtheorem{lemma}{Lemma}
\newtheorem{proposition}{Proposition}
\newtheorem{corollary}{Corollary}
\newtheorem{remark}{Remark}

\newcommand{\eps}{\varepsilon}

\title{Finite-Dimensional Quotients Obstruct Epsilon-Hypercyclicity}
\author{FA Banach run \texttt{fa\_banach\_001}, agent \texttt{agent\_lane\_08}}
\date{2026-06-19}

\begin{document}
\maketitle

\section*{Status}

This packet records a candidate partial result, likely valid.  It does not
construct a new \(\eps\)-hypercyclic non-hypercyclic operator.  Instead it
isolates a general obstruction which any solution to the open problem in
Tapia-Garcia, arXiv:2103.08075, must avoid.

\section*{Source Context}

Tapia-Garcia studies \(\eps\)-hypercyclic operators.  For a Banach space \(X\),
an operator \(T\in\mathcal L(X)\) is \(\eps\)-hypercyclic if there is a vector
\(x\in X\) such that, for every nonzero \(y\in X\), some \(n\geq 0\) satisfies
\[
  \|T^n x-y\|\leq \eps\|y\|.
\]
The introduction asks whether every separable infinite-dimensional Banach space
admits an \(\eps\)-hypercyclic operator which is not hypercyclic, with
\(\eps\in(0,1)\).

\begin{center}
\includegraphics[width=0.98\linewidth]{figures/open_problem_crop.png}
\end{center}

The source paper gives positive constructions in several classical cases,
especially when the space contains a complemented subspace of the form needed
for its operator-weighted shift construction.  The result below records a
different, negative constraint on possible constructions.

\section*{Result}

\begin{lemma}[No finite-dimensional \(\eps\)-hypercyclic operators]
\label{lem:finite-dimensional}
Let \(F\neq\{0\}\) be a finite-dimensional real or complex Banach space.  If
\(A\in\mathcal L(F)\) and \(0\leq\eps<1\), then \(A\) is not
\(\eps\)-hypercyclic.
\end{lemma}

\begin{proof}
It is enough to argue over the complexification in the real case.  Suppose that
some \(u\in F\) were an \(\eps\)-hypercyclic vector for \(A\).  Fix any
nonzero \(v\in F\).  Applying the definition to the targets \(t v\), with
\(t>0\), gives integers \(n(t)\) such that
\[
  \|A^{n(t)}u-tv\|\leq \eps t\|v\|.
\]
Thus
\[
  (1-\eps)t\|v\|\leq \|A^{n(t)}u\|
  \leq (1+\eps)t\|v\|.
\]
Letting \(t\downarrow0\) and \(t\uparrow\infty\), the orbit
\((A^n u)\) must have a subsequence converging to \(0\) and must also be
unbounded.

Let \(E=\operatorname{span}\{A^n u:n\geq0\}\), a finite-dimensional
\(A\)-invariant cyclic subspace.  Since \(A^{n_j}u\to0\) along a subsequence,
the restriction \(A_E=A|_E\) cannot have an eigenvalue \(\lambda\) with
\(|\lambda|\geq1\).  Indeed, if \(A_E^*\varphi=\lambda\varphi\) for some
nonzero \(\varphi\in E^*\), then cyclicity of \(u\) in \(E\) gives
\(\varphi(u)\neq0\); otherwise \(\varphi\) would vanish on every \(A^n u\) and
hence on all of \(E\).  But then
\[
  |\varphi(A^n u)|=|\lambda|^n|\varphi(u)|
\]
cannot tend to \(0\) along any subsequence.  Therefore all eigenvalues of
\(A_E\) have modulus strictly less than \(1\), so \(A_E^n u\to0\).  In
particular the orbit is bounded, contradicting the unboundedness forced above.
\end{proof}

\begin{proposition}[Epsilon-hypercyclicity passes to quotients with loss]
\label{prop:quotient}
Let \(q:X\to Y\) be a quotient map between Banach spaces, and suppose
\(T\in\mathcal L(X)\), \(S\in\mathcal L(Y)\) satisfy \(qT=Sq\).  If \(x\) is
an \(\eps\)-hypercyclic vector for \(T\), then \(qx\) is a \(\delta\)-hypercyclic
vector for \(S\) for every \(\delta>\eps\).
\end{proposition}

\begin{proof}
Fix \(\delta>\eps\) and \(0\neq z\in Y\).  If \(\eps=0\), choose any
preimage \(y\in X\) of \(z\).  Otherwise choose \(\eta>0\) such that
\(\eps(1+\eta)<\delta\), and then choose \(y\in X\) with \(q y=z\) and
\(\|y\|\leq(1+\eta)\|z\|\), using the quotient norm.  By
\(\eps\)-hypercyclicity there is \(n\) such that
\[
  \|T^n x-y\|\leq \eps\|y\|.
\]
Applying \(q\) gives
\[
  \|S^n qx-z\|
  =\|q(T^n x-y)\|
  \leq \eps\|y\|
  < \delta\|z\|
\]
in the case \(\eps>0\), and the same conclusion with error \(0\) in the case
\(\eps=0\).
\end{proof}

\begin{corollary}[Finite-dimensional quotient obstruction]
\label{cor:finite-quotient}
Let \(X\) be a Banach space, \(T\in\mathcal L(X)\), and \(0\leq\eps<1\).  If
\(T\) is \(\eps\)-hypercyclic, then there is no closed \(T\)-invariant subspace
\(M\subset X\) for which \(X/M\) is nonzero and finite-dimensional.
Equivalently, \(T\) has no nonzero finite-dimensional invariant quotient.
\end{corollary}

\begin{proof}
If such an \(M\) existed, \(T\) would induce an operator \(\overline T\) on the
nonzero finite-dimensional quotient \(X/M\).  Choose
\(\delta\) with \(\eps<\delta<1\).  Proposition~\ref{prop:quotient} would make
\(\overline T\) \(\delta\)-hypercyclic, contradicting
Lemma~\ref{lem:finite-dimensional}.
\end{proof}

\section*{Proof Intuition}

The point is that \(\eps<1\) is strong enough to see scale.  If an orbit
approximates every nonzero target within relative error \(\eps\), then along
the fixed ray \(\{t v:t>0\}\) its norms must come arbitrarily close to zero
and also become arbitrarily large.  Finite-dimensional linear dynamics cannot
do both.  A subsequence tending to zero rules out every spectral component of
modulus at least one on the cyclic subspace generated by the vector; hence the
orbit actually tends to zero and is bounded.

The quotient statement then says that any \(\eps\)-hypercyclic construction
with \(\eps<1\) must remain \(\delta\)-hypercyclic on every induced quotient
for all \(\delta>\eps\).  Since finite-dimensional quotients cannot have this
property, finite-dimensional marker coordinates cannot be used to force
non-hypercyclicity.

\section*{Implication for the Source Problem}

This does not answer Tapia-Garcia's question.  It does, however, remove a
natural route: decomposing a Banach space as a hyperplane plus a scalar
coordinate, making the hyperplane part hypercyclic, and using the scalar part
as a non-hypercyclic obstruction.  Such an operator would have a nonzero
finite-dimensional invariant quotient, so it cannot be \(\eps\)-hypercyclic
for any \(\eps<1\).

Thus any full construction on arbitrary separable Banach spaces must either
build the non-hypercyclic obstruction in an infinite-dimensional quotient or
subspace, or use a construction whose failure of hypercyclicity is not detected
by a finite-dimensional quotient.

\section*{Verification and Novelty Check}

The proof is non-computational.  The local run indexes were searched for the
arXiv id, epsilon-hypercyclicity terms, and finite-dimensional quotient terms;
no duplicate packet or attempt for this obstruction was found.  Web searches
for the same exact themes found the source paper and Bayart's later
sharp-threshold construction, but no exact occurrence of this
finite-dimensional quotient obstruction.

The novelty confidence is therefore moderate rather than high: the observation
is elementary and may be known folklore in linear dynamics, but it was not
visible in the bounded search and is directly useful for steering further work
on the source problem.

\begin{thebibliography}{9}

\bibitem{TapiaGarcia}
S. Tapia-Garcia,
\emph{An epsilon-hypercyclicity criterion and its application on classical
Banach spaces},
arXiv:2103.08075, 2021.

\bibitem{Bayart2023}
F. Bayart,
\emph{\(\eps\)-hypercyclic operators that are not \(\delta\)-hypercyclic for
\(\delta<\eps\)},
arXiv:2303.16677, 2023.

\end{thebibliography}

\end{document}
